% This file compiles with both LuaLaTeX and XeLaTeX
\documentclass[11pt]{article}

% Remove the "review" option to generate the final version.
\usepackage[final]{acl}
%\usepackage[review]{acl}
\usepackage{amsmath, amssymb}
\usepackage{multirow}
\usepackage[table]{xcolor}
\usepackage{threeparttable}
% This is not strictly necessary, and may be commented out,
% but it will improve the layout of the manuscript,
% and will typically save some space.
 \usepackage{microtype}
 \usepackage{placeins} % optional but helpful
% Required packages in preamble:
% \usepackage{graphicx}
\usepackage[bidi=basic]{babel}
 \usepackage{adjustbox}
 \usepackage{booktabs}
 \usepackage{multirow}
 \usepackage{rotating}
\usepackage{tikz}
 \usepackage{array}
 \usepackage{stfloats}
\usepackage{arydshln}
\usepackage{lscape}    % if rotation of whole table is ever needed
\usepackage{caption} % add this in your preamble if not already
\usepackage{expex} 
\usepackage{gb4e}
\noautomath % Optional: disables auto-math mode if causing issues
\usepackage[table]{xcolor}
\babelprovide[import]{arabic}
\babelfont[arabic]{rm}{Noto Nastaliq Urdu}

\usepackage{makecell}
\usepackage[bidi=basic]{babel}
\babelfont{rm}{TeX Gyre Termes} % similar to Times
%%% include whatever languages you need below this line
\babelprovide[import]{hindi}
\babelfont[*devanagari]{rm}{Lohit Devanagari}

\babelfont[*arabic]{rm}{Noto Sans Arabic}


%%%%%


\title{UrBLiMP:  A Benchmark for Evaluating the Linguistic Competence of Large Language Models in Urdu}
\author{
Farah Adeeba\textsuperscript{1,2}\thanks{This work was mainly conducted while the author was at UMass.},
Brian Dillon\textsuperscript{1},
Hassan Sajjad\textsuperscript{3}, Rajesh Bhatt\textsuperscript{1}
\\
\textsuperscript{1}Department of Linguistics, UMass Amherst, USA \\
\textsuperscript{2}Cluster of Excellence “The Politics of Inequality”, University of Konstanz, Germany\\
\textsuperscript{3}Faculty of Computer Science, Dalhousie University, Halifax, Canada \\
\\
%\textsuperscript{1} \texttt{\{bwdillon,bhatt\}@umass.edu} \\
Corresponding author: \texttt{farah.adeeba@uni-konstanz.de}\\
%\textsuperscript{3} \texttt{hsajjad@dal.ca}
}


\begin{document}

\maketitle


\begin{abstract}
Evaluating how large language models (LLMs) capture the grammatical structure of low-resource languages remains underexplored. This paper presents the Urdu Benchmark of Linguistic Minimal Pairs (UrBLiMP)—a diagnostic suite of 5,696 minimal pairs that contrast grammatical acceptability across ten core syntactic and morpho-syntactic phenomena in Urdu. The dataset is constructed from the Urdu Treebank and diverse text corpora, and human validation achieves a 96.1\% inter-annotator agreement, confirming its reliability. We evaluate twenty one multilingual LLMs, including LLaMA-3-70B and Gemma-3-27B-PT, and additionally assess the proprietary GPT-4o model using grammar-prompting techniques. GPT-4o (grammar-prompted) attains the highest average accuracy (97.4\%), reaching near-human performance on regular phenomena such as aspect agreement and ergativity. However, all models continue to struggle with flexible syntactic patterns like word-order variation and long-distance subject–verb agreement. UrBLiMP provides the first controlled evaluation framework for probing morpho-syntactic competence in Urdu and highlights both the progress and remaining challenges of multilingual and proprietary LLMs in low-resource settings.
\end{abstract}

\section{Introduction}

 Large Language Models (LLMs) have become crucial components of natural language processing systems, enabling a wide variety of tasks including summarization, machine translation, and dialogue generation. While evaluation efforts have largely focused on LLMs' performance in tasks requiring world knowledge and general language understanding, the extent to which these models acquire specific linguistic phenomena remains insufficiently explored—particularly for many low-resource languages.
 
 For the evaluation of linguistic knowledge of LMs, a prominent and effective methodology involves the use of minimal pairs \cite{blimp} : sequences of words that differ minimally, with only one forming an acceptable sentence (e.g., \textit{The cat \underline{sleeps}} vs. \textit{The cat \underline{sleep}}). Language Models (LMs) are evaluated on a pairwise zero-shot forced choice task, measuring whether they assign a higher probability to the acceptable versus the unacceptable sequence. 

The Benchmark of Linguistic Minimal Pairs (BLiMP) \cite{blimp}, for instance, comprises 67,000 such pairs for English, semi-automatically generated to evaluate a wide range of linguistic phenomena and assess LMs' knowledge of grammar on a large scale. Similar resources have been developed for other high-resource languages, including CLiMP \cite{climp} and SLING \cite{sling} for Chinese, TurBLiMP \cite{turBlimp} for Turkish, RuBLiMP \cite{rublimp} for Russian, JBLiMP \cite{jblimp} for Japanese, and BLiMP-NL for Dutch.

Despite the increasing success of multilingual LLMs, the question of which linguistic phenomena they can or cannot learn remains poorly understood for many languages, especially those with limited digital resources. For Urdu, comparable diagnostic resources for fine-grained linguistic evaluation are critically scarce in both coverage and granularity. While some benchmarking efforts exist \cite{urdubenchmarking}, they largely focus on functional competence. Even a dataset targeting ergativity in Hindi \cite{hindiBLiMP} has limited direct applicability to Urdu due to distinct script and vocabulary. Furthermore, MultiBLiMP \cite{multiblimp}, a multilingual extension of BLiMP, includes a small Urdu component, but with only approximately 1,000 examples focusing exclusively on subject agreement, its scope is insufficient to capture the intricate syntactic richness of the language. A comparative summary of these datasets is presented in Table \ref{tab:multilingual-datasets}.

\begin{table}[t]
\centering
\resizebox{\columnwidth}{!}{%
\renewcommand{\arraystretch}{1.5} % 1.5 times the default height
\begin{tabular}{lllll}
\hline
\rowcolor{gray!20}
           & \textbf{Language}          & \textbf{Size} & \textbf{P} & \textbf{Method} \\
\hline
BLiMP      & English                    & 67K           & 67                    & Dict \& Templates                          \\
CLiMP      & Chinese                    & 16K           & 16                    & Translation \& Templates                         \\
SLING      & Chinese                    & 38K           & 38                    & UD \& Templates                                  \\
TurBLiMP   & Turkish                    & 16K           & 16                    & Semi-automatically               \\
RuBLiMP    & Russian                    & 45K           & 45                    & Semi-automatically        \\
JBLIMP     & Japanese                   & 331           & 39                    & Extracted from articles                         \\
BLiMP-NL   & Dutch                      & 8.4K          & 84                    &  Semi-automatically                    \\
MultiBLiMP & 101 Languages              & 128K          & 2                     & UD                          \\
BHS           & \makecell[tl]{ Basque, Hindi,\\ Swahili}     & 300           & 3                     & Dictionary \& Templates                         \\
\textbf{UrBLiMP} & Urdu                 & 5696            & 19                    & Text corpus, rules                              \\
\hline
\end{tabular}%
}
\caption{Comparison of existing minimal pair datasets across languages. P is number of paradigms,UD: Universal Dependencies, BHS:  work by \cite{hindiBLiMP}}
\label{tab:multilingual-datasets}
\end{table}

This work aims to bridge this significant gap in Urdu by introducing \textbf{UrBLiMP}—the \textit{Urdu Benchmark of Linguistic Minimal Pairs}—a syntactically informed diagnostic dataset for evaluating LLMs' grammatical competence in Urdu. UrBLiMP comprises 5,696 minimal pairs spanning ten core syntactic phenomena. These were constructed through a hybrid methodology combining treebank-derived templates with surface pattern matching applied to diverse raw corpora, followed by rigorous manual verification.

By providing this benchmark, we offer a linguistically grounded tool for identifying systematic weaknesses in LLMs’ understanding of Urdu grammar. UrBLiMP thus lays the foundation for advancing both the evaluation and improvement of multilingual models in the context of low-resource, morphologically rich languages like Urdu.

\section{UrBLiMP}
The process of dataset creation is outlined in this section. The selected linguistic phenomena and corresponding paradigms are shown in Table~\ref{tab:urdu-phenomena}. Minimal pairs were constructed using both the Urdu Treebank and surface-level patterns applied to raw text. For each phenomenon, transformation rules were formulated to produce minimal pairs differing on a single grammatical property. The following subsections describe the selected linguistic phenomenon, minimal pairs generation and transformation procedures in more detail.
\FloatBarrier  % optional: prevent floats moving past this point
\begin{table*}[ht!]
\resizebox{\textwidth}{!}{
    \begin{tabular}{llp{9cm}p{9cm}}
    \hline
    \rowcolor{gray!20}
    \textbf{Phenomenon} & \textbf{N} & \textbf{Grammatical Sentence} & \textbf{Ungrammatical Sentence} \\
    \hline

    Aspect & 1 
    & \parbox[t]{9cm}{
       \foreignlanguage{arabic}{ہاں وہ گورنر سے ملتا \underline{رہا} تھا۔} \\
\textit{        h\~{a} \quad voh \quad gavarnar \quad se \quad milta: \quad raha: \quad tha:} \\
        Yes \quad he \quad governor \quad from \quad meet-\textsc{impf} \quad \textsc{prog.part} \quad was \\
        `Yes, he had been meeting the governor.' \\
    }
    & \parbox[t]{9cm}{  
    \foreignlanguage{arabic}{ہاں وہ گورنر سے ملتا \underline{چکا} تھا۔} \\
\textit{    h\~{a} \quad voh \quad gavarnar \quad se \quad milta: \quad cuka: \quad tha:} \\
    Yes \quad he \quad governor \quad from \quad meet-\textsc{impf} \quad \textsc{perf.part} \quad was \\
    `Yes, he had already met the governor.' \\} \\ \hline

    Dative Object & 2 & \parbox[t]{9cm}{  \foreignlanguage{arabic}{پہرے‌دار نے تمام ماجرا راجا \underline{کو} سنایا۔} \\
   pahreda:r \quad ne \quad tama:m \quad ma:jra: \quad ra:ja: \quad ko \quad suna:ya: \\
   guard \quad \textsc{erg} \quad whole \quad incident \quad Raja \quad \textsc{dat} \quad narrated.\textsc{pst} \\
   `The guard narrated the whole incident to Raja.' \\ }

 &  \parbox[t]{9cm}{   \foreignlanguage{arabic}{پہرے‌دار نے تمام ماجرا راجا سنایا۔} \\
   pahreda:r \quad ne \quad tama:m \quad ma:jra: \quad ra:ja: \quad suna:ya: \\
   guard \quad \textsc{erg} \quad whole \quad incident \quad Raja \quad narrated.\textsc{pst} \\
   `The guard narrated the whole incident to Raja.'  }
 \\ \hline

 Ergativity & 3 & \parbox[t]{9cm}{  \foreignlanguage{arabic}{پہرے‌دار نے تمام ماجرا راجا کو \underline{سنایا}} \\
   pahreda:r \quad ne \quad tama:m \quad ma:jra: \quad ra:ja: \quad ko \quad suna:ya: \\
   guard \quad \textsc{erg} \quad whole \quad incident \quad Raja \quad \textsc{dat} \quad narrated.\textsc{pst} \\
   `The guard narrated the whole incident to Raja.' \\ }
    &  \parbox[t]{9cm}{  \foreignlanguage{arabic}{پہرے‌دار نے تمام ماجرا راجا کو \underline{سناتا}} \\
   pahreda:r \quad ne \quad tama:m \quad ma:jra: \quad ra:ja: \quad ko \quad suna:ta: \\
   guard \quad \textsc{erg} \quad whole \quad incident \quad Raja \quad \textsc{dat} \quad narrates.\textsc{pres} \\
   `The guard narrated the whole incident to Raja.' \\ }\\ \hline

   Experiencer Subject & 1 & \parbox[t]{9cm}{

\foreignlanguage{arabic}{ \underline{مجھے} یہ بلاگ پوسٹ پسند آیا۔} \\
\textit{mujhe yeh blog post pasand a:ya:} \\
\textit{to.me this blog post liking came} \\
I liked this blog post.

} & \parbox[t]{9cm}{\foreignlanguage{arabic}{ \underline{میں} یہ بلاگ پوسٹ پسند آیا۔ } \\
\textit{maĩ yeh blog post pasand a:ya:} \\
\textit{I this blog post liking came} \\
I liked this blog post.} \\ \hline

Honorific & 1 & 
\parbox[t]{9cm}{\foreignlanguage{arabic}{ بہن جی مجھ سے ناراض \underline{تھیں} }\\
\textit{behan ji mujh se nara:z th\~{i}} \\
\textit{sister from.me upset was (hon.)} \\
Sister was upset with me.}
& 
\parbox[t]{9cm}{\foreignlanguage{arabic}{ بہن جی مجھ سے ناراض \underline{تھی}} \\
\textit{behan ji mujh se nara:z thi} \\
\textit{sister from.me upset was} \\
Sister was upset with me.} \\
\hline
N-J Agr & 1 & 
\parbox[t]{9cm}{\foreignlanguage{arabic}{ ہر مجسمہ بیس فٹ \underline{اونچا} ہے ۔} \\
\textit{har mujasma bi:s fit o:nca: hai} \\
\textit{every statue twenty feet tall is} \\
Every statue is twenty feet tall.}
& 
\parbox[t]{9cm}{\foreignlanguage{arabic}{ ہر مجسمہ بیس فٹ \underline{اونچی} ہے ۔} \\
\textit{har mujasma bi:s fit onci: hai} \\
\textit{every statue twenty feet tall.F is} \\
Every statue is twenty feet tall.} \\
\hline
Obliqueness& 5 &\parbox[t]{9cm}{\foreignlanguage{arabic}{\underline{بڑے} صاحب نے کہا}\\
\textit{baRe sahib ne kaha:}\\
\textit {big.\textsc{obl.m.sg} boss.\textsc{nom} \textsc{erg} say.\textsc{pfv.m.sg} }\\
\textit{big boss said}

} & 
\parbox[t]{9cm}{
\foreignlanguage{arabic}{\underline{بڑا} صاحب نے کہا}\\
\textit{baRa: sahib ne kaha:}\\
\textit{ big.\textsc{nom.m.sg} boss.\textsc{nom} \textsc{erg} say.\textsc{pfv.m.sg}}
\\
\textit{big boss said}} 
\\

\hline
Participial Relatives & 1 &
\parbox[t]{9cm}{
\foreignlanguage{arabic}{آسانی سے \underline{کھولا} گیا دروازہ بند کیا جا سکتا ہے} \\
\textit{a:sani: se khola: gaya: darwaza: band kiya: ja: sakta: hai} \\
easy-\textsc{adv} from open-\textsc{ptcp} go-\textsc{ptcp.m.sg} door close do-\textsc{perf.m.sg} go can be \\
‘The door that was easily opened can be closed.’
}
&
\parbox[t]{9cm}{
\foreignlanguage{arabic}{آسانی سے \underline{کھولتا} گیا دروازہ بند کیا جا سکتا ہے} \\
\textit{a:sani: se kholta: gaya: darwaza: band kiya: ja: sakta: hai} \\
easy-\textsc{adv} from open-\textsc{ipfv.m.sg} go-\textsc{ptcp.m.sg} door close do-\textsc{perf.m.sg} go can be \\
‘The door that kept opening easily can be closed.’
} \\
\hline
       Subj-Verb Agr & 3 & 
  \parbox[t]{9cm}{
\foreignlanguage{arabic}{اور اس کا باپ تو شاید پاگل ہو \underline{جاتا} ۔} \\
\textit{aur us ka ba:p to shayad pagal ho jata:} \\
and 3\textsc{sg.m} \textsc{gen.m} father \textsc{focus} perhaps mad become \textsc{cond.m.sg} \\
‘And his father might have gone mad.’
}
&
\parbox[t]{9cm}{
\foreignlanguage{arabic}{اور اس کا باپ تو شاید پاگل ہو \underline{جاتی} ۔} \\
\textit{aur us ka ba:p to shayad pagal ho jati:} \\
and 3\textsc{sg.m} \textsc{gen.m} father \textsc{focus} perhaps mad become \textsc{cond.f.sg} \\
‘And his father (incorrectly) might have gone mad (feminine verb).’
} \\ \hline

Order Variation& 1 &
\parbox[t]{9cm}{
\foreignlanguage{arabic}{جو \underline{بات} ہے وہ بولو۔} \\
\textit{jo ba:t hai vo bolo} \\
what matter is that say.\textsc{imp.pl} \\
‘Say what the matter is.’
}
&
\parbox[t]{9cm}{
\foreignlanguage{arabic}{جو ہے وہ بولو بات ۔} \\
\textit{jo hai vo bolo ba:t} \\
what is that say matter \\
‘Say what is that matter.’ (incorrect order)
} \\ \hline

    \end{tabular}
    }
    \caption{Ten Urdu linguistic phenomena covered in this study, with grammatical and ungrammatical sentence examples. Minimal contrasts are emphasized. The second line of each example shows transliterated Urdu, and the third line provides an English translation. \textit{N} indicates the number of paradigms within each phenomenon.}
    \label{tab:urdu-phenomena}

\end{table*}
\subsection{Linguistic Phenomenon}
UrBLiMP\footnote{UrBLIMP data can be found at: \href{https://github.com/farahadeeba/UrBLiMP}{https://github.com/farahadeeba/UrBLiMP}} covers ten major linguistic phenomena of Urdu, as summarized in Table \ref{tab:urdu-phenomena}.  

\paragraph{Selection criteria.}
The choice was guided by both linguistic significance and typological relevance to Urdu. Specifically, we aimed to capture phenomena that are:

\begin{itemize}
    \item \textbf{Universally relevant:} Core syntactic dependencies that are broadly attested across languages and central to sentence well-formedness (e.g., Subject–Verb Agreement).
    \item \textbf{Typologically distinctive:} Properties that are characteristic of Indo-Aryan languages and particularly salient in Urdu (e.g., ergative alignment in the perfective, honorific agreement).
    \item \textbf{Morphosyntactically rich and diagnostically clear:} Domains that permit precise minimal contrasts with surface-diagnosable cues (e.g., Noun–Adjective Agreement, Obliqueness and its interaction with postpositions).
\end{itemize}
Several paradigms are designed around a linguistic phenomenon to check the LMs robustness.

\textbf{Aspect Agreement} phenomena capture the restrictions governing co-occurrence of aspectual markers in Urdu. Specifically, while the habitual participle {\em -ta:}  ( \foreignlanguage{arabic}{تا/-تی/-تے- }) may combine with the progressive marker {\em raha:}  ( \foreignlanguage{arabic}{رہا، رہی ، رہے }) to express ongoing habitual actions, its combination with the perfective aspect verb {\em cuka:}   ( \foreignlanguage{arabic}{چکا})  is ungrammatical. The mismatch arises because the habitual aspect denotes iterative or unbounded events, whereas the perfective aspect verb implies a bounded and completed action \cite{butt2008}, leading to a semantic conflict.

\textbf{Dative Object} phenomenon investigate the ungrammaticality that results from omitting the dative postposition {\em ko} ( \foreignlanguage{arabic}{ کو}) when it is required to mark direct objects in Urdu. Two paradigms were considered: 1) the direct object is a proper noun and 2) the direct object 
is a pronoun. 

\textbf{Ergativity} in Urdu is a split system primarily governed by aspect. In perfective transitive constructions, the subject takes ergative marking {\em ne} ( \foreignlanguage{arabic}{نے}) and the verb agrees with the direct object. In non-perfective aspects (e.g., habitual or progressive), the subject is nominative and the verb agrees with it. To assess model sensitivity to this system, minimal pairs across three paradigms are created:
\begin{enumerate}
    \item \textit{Aspect Sensitivity:} Ungrammaticality was introduced by replacing the perfective verb with a habitual verb while retaining the ergative subject. This results in a mismatch between aspect and subject marking.
    
    \item \textit{Verb-Object Agreement:} In perfective transitive clauses, the verb should agree in gender with the direct object rather than the ergative-marked subject. Ungrammatical variants were created by forcing verb agreement with the subject instead.
    
  \item \textit{Differential Object Marking (DOM) in Ergative Constructions:}
In Urdu ergative constructions (i.e., perfective transitive clauses with the subject marked by \textit{ne}), when the direct object is marked with \textit{ko}, the verb typically defaults to masculine singular form, regardless of the object’s gender or number. To test sensitivity to this default agreement rule, ungrammatical sentence variants were constructed by replacing the expected masculine singular verb with a feminine form that agrees with the feminine object.


\end{enumerate}

\textbf{Experiencer Subjects.} In Urdu, experiencer-subject constructions typically involve a subject marked in the dative case, often with the postposition \textit{ko}. However, when the subject is a pronoun, the dative case is realized through dedicated oblique pronominal forms (e.g., \textit{mujhe}, \textit{tujhe}), and the postposition \textit{ko} is not overtly used. 

In this paradigm, we focus on constructions where the experiencer is a pronominal subject in its dative (oblique) form . Ungrammatical variants were generated by replacing these forms with nominative pronouns (e.g., \textit{meN}) or incorrect oblique forms that lack proper dative alignment (e.g., \textit{us} without \textit{ko}).

\textbf{Honorific.} Urdu marks respect or politeness by enforcing plural agreement on the verb when addressing or referring to someone honorifically. This applies even when the referent is singular, especially in constructions involving titles or respectful nouns such as (ji:, muhtram, muhtarma:, ja:n)\foreignlanguage{arabic}{محترم}, \foreignlanguage{arabic}{محترمہ}, \foreignlanguage{arabic}{جی}, or \foreignlanguage{arabic}{جان}. To evaluate this phenomenon, ungrammatical variants were created by replacing the required plural verb forms with singular forms, thereby violating the honorific agreement pattern. 

\textbf{Obliqueness} phenomenon tests 
%whether 
models 
%can detect 
for 
case-driven contrasts where oblique forms are required due to syntactic context. We construct ungrammatical variants 
%were constructed 
by replacing expected oblique forms with direct (non-oblique) forms across various categories, including masculine singular nouns, plural nouns, adjectives, verbs, and pronouns.


\textbf{Participial Relatives} phenomenon involves the use of perfective participles in relative clauses that modify a noun e.g. \textit{ khola gaya darvaza:(opened door)}. In Urdu, such constructions typically require a passive participial verb form to ensure grammaticality. To test this, ungrammatical variants were generated by replacing the participial verb with an imperfective or active form, disrupting the required agreement and aspectual constraints within the relative clause. 

\textbf{Subject-Verb Agreement} phenomenon is tested across three paradigms: number, gender, and person agreement. Ungrammatical sentence variants were generated by violating the expected agreement between the subject and the verb in each paradigm.




\textbf{Word Order Variation}
Urdu generally allows relatively flexible word order; however, certain positions are preferred for clarity and naturalness, especially in relative clauses. In this paradigm, ungrammatical  or less acceptable variants were created by altering the canonical word order, such as shifting the noun out of a relative clause to the clause-final position, which disrupts the natural syntactic structure and interpretability of the sentence.

\subsection{Minimal Pairs Generation}
 Minimal pairs were generated using naturally occurring sentences extracted from Urdu texts. Two primary resources were utilized for this purpose: the Urdu Treebank and in-house Urdu Corpus.
\subsubsection{Urdu Treebank}
To extract sentences exhibiting the targeted linguistic phenomena, the Urdu Treebank \cite{cletreebank} was utilized. This resource comprises 7,854 Urdu sentences annotated in the Penn Treebank style. A subset containing relevant syntactic structures was selected using pattern-based search. 

For example, to extract sentences illustrating the Aspect Agreement phenomenon, the syntactic pattern \textit{(VC (VBF) (AUXP))} was used. From the matching sentences, those containing a VBF verb with a habitual participle suffix (e.g., -ta) were retained. These extracted sentences were then manipulated to construct minimal pairs: the auxiliary phrase (AUXP) — such as \textit{raha:, rahi:} — was replaced with \textit{cuka: or cuki:} to generate ungrammatical variants. 
\subsubsection{Urdu Corpus}
Since a sufficient number of required sentences could not be obtained from the Urdu Treebank, additional data were extracted from an in-house collected corpus of Urdu text. This corpus comprises approximately 735 million tokens and includes diverse web-based content from sources such as Common Crawl, Twitter, and news articles.

For each paradigm, specific regular expressions were crafted to extract potentially relevant sentences. For example, to retrieve candidate sentences for the word order variation phenomenon, the regular expression 
\textit{\^{}\foreignlanguage{arabic}{جو}\textbackslash b.*\textbackslash b \foreignlanguage{arabic}{ وہ }\textbackslash b}
was used. The extracted sentences were then manually reviewed, and only those that conformed to the targeted linguistic paradigm were retained. Furthermore, data cleaning was performed to correct typographical errors and remove unnecessary space insertions or deletions. Finally, minimal pairs were constructed from the selected sentences by systematically introducing controlled variations.

\subsection{Human Evaluation}
To assess the quality of the minimal pairs generated, two rounds of human validation were performed using the PCIbex platform \cite{pcibex}. A total of 17 native Urdu speakers participated in the evaluation. The group included six males and eleven females, with ages ranging from 23 to 46 years. Among the participants, four were linguists, while the remaining evaluators had at least a high school education, including three with Ph.D.-level qualifications.

In the first round, all annotators underwent a brief training phase in which they at least annotated 20 pairs of demonstration sentences. During both the training and the main task, annotators were instructed to make pairwise acceptability judgments: given a pair of sentences, indicate which sentence is more acceptable. This familiarization step helped ensure a consistent understanding of the evaluation interface, task, and decision criterion.

In the second round, each annotator labeled approximately 190 pairs, covering about 10 pairs from each linguistic paradigm. For each pair, annotators selected the more acceptable sentence; absolute grammaticality labels were not required. The sentences of different paradigms were randomly shuffled. To ensure the reliability of judgments, each sentence pair was annotated by at least three different evaluators. The complete annotation task took approximately one hour per annotator. The final raw human accuracy mean overall paradigms is 96.10\%. The inter-annotator agreement as measured by Fleiss’ kappa is 0.89, indicating almost perfect agreement according to \cite{agreement}. The average acceptability accuracy per paradigm is shown in Table \ref{tab:human_llm_subset_appendix}.
%To assess the quality of the minimal pairs generated, two rounds of human validation were performed using the PCIbex platform \cite{pcibex}. A total of 17 native Urdu speakers participated in the evaluation. The group included six males and eleven females, with ages ranging from 23 to 46 years. Among the participants, four were linguists, while the remaining evaluators had at least a high school education, including three with Ph.D.-level qualifications.

%In the first round, all annotators underwent a brief training phase in which they at least annotated 20 pairs of demonstration sentences. This familiarization step helped ensure a consistent understanding of the evaluation interface and task.

%In the second round, each annotator labeled approximately 190 pairs, covering about 10 pairs from each linguistic paradigm. The sentences of different paradigms were randomly shuffled. To ensure the reliability of judgments, each sentence pair was annotated by at least three different evaluators. The complete annotation task took approximately one hour per annotator. The final raw human accuracy mean overall paradigms is 96.10\%. The inter-annotator agreement as measured by Fleiss’ kappa is 0.89, indicating almost perfect agreement according to \cite{agreement}. The average acceptability accuracy per paradigm is shown in Table \ref{tab:human_llm_subset_appendix}.

\section{Experimental Setup}
This section presents the experimental setup used to evaluate the syntactic understanding of various multilingual and Urdu language models. We detail the selected models and describe the accuracy-based evaluation metric applied across linguistic phenomena.

\subsection{Evaluation Models}
For evaluation, we employ a diverse set of 20 multilingual 
language models. 

These include encoder-only models such as multilingual \textbf{BERT} \cite{bert} , encoder–decoder models like \textbf{mT5} \cite{mt5}, and decoder-only models including \textbf{LLaMA3} \cite{llama3} and \textbf{Gemma} \cite{gemma_2025}. The Gemma models are particularly noteworthy for being trained on a balanced multilingual corpus encompassing 140 languages. 

We also consider \textbf{BLOOMZ} \cite{bloomz}, \textbf{DeepSeek} \cite{deepseekai} and 
%models in our evaluation of UrBLiMP. In addition, models from 
the \textbf{Granite} series \cite{granite} models. % were evaluated.
Moreover, a notable addition is \textbf{Alif-1.0-8B-Instruct} \cite{alif}, a publicly available, continually pre-trained Urdu model based on \texttt{unsloth/Meta-Llama-3.1-8B}.

For closed-weights model, we evaluated GPT-4.0 via the provider API using prompt-based acceptability tests. We followed the \emph{explain-then-process} (``grammar prompting'') paradigm proposed by Scheinberg et al., which first elicits a concise grammatical rule and then asks the model to apply that rule; we implemented three prompting regimes: (i) \textbf{base} (direct acceptability or two-choice template), (ii) \textbf{GP} (grammar prompting), and (iii) \textbf{GP+CoT} (grammar prompting augmented with an explicit chain-of-thought). The GP design is inspired by Scheinberg et al. \cite{scheinberg} and details of the exact GPT-4.0 prompt templates and extraction procedure are provided in Appendix \ref{appendix:prompts}.

\subsection{Evaluation Measure}
Following prior work \cite{sling, multiblimp}, we perform evaluation 
%is performed 
at the sentence level, rather than at the position of the key linguistic items, which was the focus of earlier studies. Consistent with \citep{sling}, we compute perplexity \cite{perplexity} for causal language models and pseudo-perplexity for masked language models over each linguistic pair.

To evaluate a language model, we use \textbf{accuracy} as the primary metric. Accuracy is defined as the proportion of minimal pairs in which the model assigns a lower (pseudo-)perplexity to the grammatically acceptable sentence.

Formally, let $N$ be the total number of sentence pairs $(s_{\text{good}}, s_{\text{bad}})$, where $s_{\text{good}}$ is the acceptable sentence and $s_{\text{bad}}$ is the unacceptable counterpart. Let $\text{ppl}(s)$ denote the (pseudo-)perplexity of sentence $s$. Then the accuracy is computed as:

\begin{equation*}
\text{Accuracy} = \frac{1}{N} \sum_{i=1}^{N} \mathbb{I} \left[ \text{ppl}(s_{\text{good}}^i) < \text{ppl}(s_{\text{bad}}^i) \right]
\end{equation*}

where $\mathbb{I}[\cdot]$ is the indicator function that returns 1 if the condition holds and 0 otherwise.
\section{Results and Analysis}
Table \ref{tab:syntactic-eval} presents the results of the language models (LMs) and human performance on each syntactic phenomenon. Overall, LLaMA-3-70B achieves the highest average performance across most paradigms. However, there is no single LM that consistently outperforms all others across all categories. In several instances, other LMs surpass LLaMA-3-70B on specific syntactic phenomena.

A comparison between multilingual models and the continual-pretrained model on Urdu (i.e., Alif-LLaMA-8B) reveals that Alif achieves a higher average accuracy than its base model, LLaMA-3-8B, based on its performance across the linguistic phenomena in Table~\ref{tab:syntactic-eval}. While Alif does not consistently outperform all models on individual phenomena, it shows a notable improvement on the Word Order Variation phenomenon and achieves comparable performance on the Experiencer-Subject construction. These findings underscore the benefits of continued pretraining on Urdu data.



\begin{table*}[ht!]
\centering
\begin{threeparttable}

\resizebox{\textwidth}{!}{
\renewcommand{\arraystretch}{1.5} % 1.5 times the default height
\begin{tabular}{lrrrrrrrrrrr}
 & \multicolumn{1}{l}{\rotatebox{70}{Aspect Agr}} & \multicolumn{1}{l}{\rotatebox{70}{Dative-Obj} }& \multicolumn{1}{l}{\rotatebox{70}{Ergativity}} & \multicolumn{1}{l}{\rotatebox{70}{Exper-Sub}} & \multicolumn{1}{l}{\rotatebox{70}{Honorific}} & \multicolumn{1}{l}{\rotatebox{70}{N-J Agr}} & \multicolumn{1}{l}{\rotatebox{70}{Oblique}} & \multicolumn{1}{l}{\rotatebox{70}{PR}} & \multicolumn{1}{l}{\rotatebox{70}{S-Verb Agr}} & \multicolumn{1}{l}{\rotatebox{70}{ order vartion}} & \multicolumn{1}{l}{\rotatebox{70}{\textbf{Average}}} \\
\hline
gemma-3-1b-pt & 100.00  & 98.62  & 87.36  & 79.01  & 73.20  & 93.00  & 87.34  & 79.07  & 87.60  & 82.18  & 86.74  \\
gemma-3-4b-pt & 100.00  & 98.24  & 90.15  & 84.20  & 73.86  & 93.00  & 85.58  & 81.06  & 88.91  & 86.14  & 88.11  \\
gemma-3-12b-pt & 100.00  & 98.24  & 93.41  & 94.07  & 84.97  & 96.50  & 90.20  & 85.71  & 92.58  & 91.09  & \textbf{92.68 } \\
gemma-3-27b-pt & 100.00  & \textbf{99.62 } & 94.94  & 84.20  & 84.31  & \textbf{97.00 } & \textbf{92.75 } & 85.38  & \textbf{94.89 } & 91.09  & \textbf{92.42 } \\
gemma-3-1b-it & 99.00  & 89.81  & 70.02  & 62.72  & 55.56  & 76.50  & 70.07  & 57.48  & 70.44  & 63.37  & 71.50  \\
gemma-3-4b-it & 99.75  & 92.65  & 86.62  & 56.54  & 74.51  & 77.00  & 76.78  & 74.42  & 82.19  & 68.32  & 78.88  \\
gemma-3-12b-it & 100.00  & 95.16  & 91.72  & 76.05  & 79.74  & 89.50  & 79.83  & 78.74  & 89.99  & 78.22  & 85.89  \\
gemma-3-27b-it & 100.00  & 94.40  & 93.98  & 56.05  & 84.97  & 92.50  & 85.05  & 85.38  & 92.55  & 83.17  & 86.81  \\
Llama-3-8B & 100.00  & 92.60  & 91.45  & \textbf{98.27 } & 93.46  & 89.50  & 85.14  & 90.03  & 76.60  & 90.10  & 90.71  \\
Llama-3-70B & 100.00  & 96.43  & \textbf{96.38 } & 95.80  & 98.04  & 94.00  & \textbf{92.77 } & \textbf{92.69 } & 87.13  & 94.06  & \textbf{94.73 } \\
Alif-1.0-8B-Instruct & 100.00  & 91.26  & 93.61  & 98.02  & 98.04  & 89.00  & 90.73  & 91.69  & 83.13  & \textbf{97.03 } & \textbf{93.25 } \\
granite-3.3-8b-base & 97.51  & 88.83  & 88.67  & 91.11  & 90.20  & 85.00  & 84.83  & 82.72  & 72.48  & 93.07  & 87.44  \\
granite-3.3-2b-base & 99.25  & 89.59  & 82.07  & 77.04  & 83.66  & 83.00  & 82.25  & 80.07  & 64.78  & 85.15  & 82.69  \\
granite-3.3-8b-instruct & 98.75  & 89.21  & 84.14  & 60.74  & 94.77  & 80.50  & 77.92  & 74.09  & 71.06  & 91.09  & 82.23  \\
granite-3.3-2b-instruct & 99.00  & 80.53  & 80.18  & 92.10  & 88.89  & 74.00  & 74.04  & 78.41  & 61.07  & 76.24  & 80.45  \\
DeepSeek\tnote{1} & 100.00  & 88.20  & 93.46  & 96.30  & \textbf{99.35 } & 90.50  & 91.66  & 88.70  & 86.33  & 90.10  & 92.46  \\
bloomz-7b1-p3 & 99.25  & 96.28  & 90.85  & 67.65  & 84.31  & 87.50  & 84.91  & 81.06  & 84.58  & 85.15  & 86.15  \\
gpt-4o(base) &98.00 &94.28 & 93.41& 98.27&86.93 &\textbf{99.00}  &96.22  & 96.01&96.70 & 100 & 95.69\\

gpt-4o(GP) &99.00 &97.54 & 93.66&99.01 &98.69 & 98.00&\textbf{97.60} &\textbf{97.67}& \textbf{98.06} & \textbf{100} & \textbf{97.40} \\

gpt-4o(GP+CoT) &96.26 & 95.09&93.69 &\textbf{99.26} &97.39 &93.50&92.58&97.34 & 91.29& 99.01 &94.25 \\

Bert\tnote{2} & 70.82 & 75.44 & 79.91 & 82.02 & 80.97 & 81.36 & 80.95 & 79.50 & 69.73 & 90.10 & 78.77 \\
mt5-small & 27.43 & 29.77 & 54.75 & 60.62 & 36.60 & 46.50 & 60.67 & 79.40 & 50.79 & 60.40 & 50.69 \\
mt5-large & 71.82& 58.12& 58.14& 79.75& 51.50& 56.48& 63.56& 73.52& 58.30& 62.38& 63.36\\


\rowcolor{gray!20}
\textbf{Human} & \textbf{100} & \textbf{94.27} & \textbf{97.71} & \textbf{93.89} & \textbf{94.17} & \textbf{97.29} & \textbf{95.46} & \textbf{95.63} & \textbf{95.68} & \textbf{98.75} \\

\bottomrule
\end{tabular}
}
\caption{Percentage accuracy of various models on  different syntactic phenomena. Average accuracy accross all phenomena, it = Instruction Tuned, pt = Pretrained, PR = Participial Relatives}

\begin{tablenotes}
\item[1] DeepSeek-R1-Distill-Llama-70B
\item[2] bert-base-multilingual-cased
\end{tablenotes}
\label{tab:syntactic-eval}
\end{threeparttable}
\end{table*}


\subsection{Long Distance Agreement}
%
%For example, 
In Urdu, gender agreement between the subject and verb is crucial. 
%However, as the syntactic distance between these elements grows, errors in maintaining this agreement are often produced by the language model. 
For example in (\ref{ex:gender1}) verb is separated from  the noun. This indicates that the capability of the model to capture and utilize long-range syntactic relationships is limited, which negatively impacts its overall performance on such linguistic phenomena.

We observed that the language models struggle particularly with long-distance agreement. Specifically, in cases such as gender paradigm in phenomenon of subject-verb agreement, when the distance between the subject and the verb increases, the ability of the model to correctly predict gender agreement significantly deteriorates.

For example (\ref{ex:gender2}) is marked as ungrammatical by most of the LMs even with the LlaMA-70B model because there is long distance between  subject \foreignlanguage{arabic}{استانی} (teacher.F) and  verb \foreignlanguage{arabic}{سوتی}(sleep.F). 

\begin{exe}

\ex\label{ex:gender1}
\foreignlanguage{arabic}{فیصل اپنے سکول سے خوشی خوشی گھر واپس آیا}\\Faisal apne sku:l sy Kushi Kushi gher aya: \\Faisal.M his.POSS school from happily home came-Perf.M
\glt Faisal came home happily from his school
\ex \label{ex:gender2}
\foreignlanguage{arabic}{استانی سبق پڑھا کر آرام سے کرسی پر بیٹھ کر سوتی۔
}\\
Ustani: sabq peRha: kar a:ra:m se kursi: par beTh kar soti:\\
teacher lesson teach-Perf.F do-Perf.F comfortably chair on sit-Perf.F sleep-Perf.F
\glt 'After teaching the lesson, the teacher would sit comfortably on a chair and sleep.'

\end{exe}

\begin{figure}[t]
  \centering
  \includegraphics[width=\columnwidth]{latex/images/gender.png}
  \caption{Best-performing model from each architecture family on the subject-Verb Agreement in gender paradigm.}
  \label{fig:gender_best_family}
\end{figure}
Figure~\ref{fig:gender_best_family} presents a comparison of the best-performing model from each architecture family evaluated on the gender paradigm of Subj-Verb phenomenon. The \texttt{gemma-3-27b-pt} model achieved the highest accuracy , substantially outperforming others, including LLaMA-3-70B  , DeepSeek’s LLaMA-70B  , \texttt{Alif-8B-Instruct} and \texttt{Bloomz-7b1-p3}.

In contrast, smaller encoder-based models such as \texttt{bert-base-multilingual-cased} and \texttt{mt5-large} achieved considerably lower accuracy (67.42\% and 57.30\%, respectively), indicating the limitations of encoder-only or smaller multilingual models for fine-grained gender-related linguistic understanding. 

\subsection{Phenomenon Specific Results}
A linguistic breakdown of the evaluation results reveals varying degrees of model competence across different grammatical phenomena in Urdu. Some constructions—particularly those that are categorically distinct in surface form—proved relatively straightforward for both humans and language models, whereas others remained more challenging, especially for smaller models.

Aspect Agreement emerged as one of the simplest phenomena for both human annotators and LLMs. Even smaller-scale models achieved near-perfect accuracy on this task. A likely explanation is that in Urdu, habitual and perfective aspects are rarely co-expressed within the same clause, resulting in clear distributional patterns. This distinct separation made the classification task more deterministic and less ambiguous for models.

In the case of Dative Object marking, both noun and pronoun contexts yielded accuracies above 92\% across most models (see  Table \ref{tab:pretrain-paradigm}). However, models such as those in the Granite series displayed slightly reduced performance for dative nouns, achieving accuracies around 88\%. Pronouns, on the other hand, were processed with higher precision, likely due to their more regular and less variable surface forms.

With respect to Ergativity, models generally succeeded in detecting aspect-sensitive agreement patterns, particularly the requirement for perfective verb forms in ergative constructions. Nonetheless, performance dropped considerably in related phenomena such as Object-Verb Agreement and Differential Object Marking (DOM). In both these cases, most models exhibited accuracy scores below 90\%, suggesting that these constructions involve more nuanced syntactic and semantic dependencies, which continue to challenge current LLMs.

In Obliqueness, models struggled particularly with adjective and masculine singular nouns even when its very local. The oblique case marking in Urdu—especially when realized through subtle morphological alternations—proved difficult for several models to capture consistently. Only \texttt{LlaMA-70B and DeepSeek-70B} accuracy is above 90\%. This indicates an ongoing limitation in how models process inflectional morphology in morphologically rich languages like Urdu.

In the case of subject–verb agreement, the highest accuracy of 94.89\% is achieved. The models tend to struggle with gender agreement, where the maximum accuracy reaches 92.51\%, particularly when the subject and verb are separated by greater syntactic distance.
To examine whether these rankings generalize across benchmarks, we also evaluated the same models on the Urdu subject–verb agreement subset of MultiBLiMP (550 examples); the detailed results are presented in Appendix~\ref{appendix:multiblimp}. The comparison confirms that model behavior is consistent across benchmarks: accuracy is systematically lower on UrBLiMP, reflecting its greater difficulty, yet the Gemma-3-27B-PT model remains the top performer in both evaluations.  

\subsection{Model Size}
\begin{table}[!t]


\resizebox{\columnwidth}{!}{
\begin{tabular}{lllll}
\toprule
\textbf{LM 1} & \textbf{LM 2} & \textbf{p-raw} & \textbf{p-holm} & \textbf{Sig-holm} \\
\midrule
gemma-3-1b-pt & 4b-pt & 0.078 & 0.156 & No \\
gemma-3-1b-pt & 12b-pt & 0.007 & 0.039& Yes \\
gemma-3-1b-pt & 27b-pt & 0.003 & 0.023 & Yes \\
gemma-3-4b-pt & 12b-pt & 0.007 & 0.039 & Yes \\
gemma-3-4b-pt & 27b-pt & 0.007 & 0.039 & Yes \\
gemma-3-12b-pt & 27b-pt & 0.460 & 0.460 & No \\
\bottomrule
\end{tabular}
}
\centering
\caption{Pairwise Wilcoxon signed-rank test results for Gemma-PT models on linguistic phenomena. Holm-corrected p-values are reported.}
\label{tab:gemma_pairwise}
\end{table}

We observed that model performance generally improved with increasing parameter size within the Gemma series. Notable accuracy gains were achieved between the 1B and 27B parameter models, while there is minimal decline in accuracy from 12B to 27B. The statistical significance of is gain was confirmed through a two-tailed Wilcoxon signed-rank test. Table~\ref{tab:gemma_pairwise} presents the results among the \texttt{gemma-PT} models to assess whether the differences in performance are statistically significant. Table~\ref{tab:gemma_pairwise} shows the difference between \texttt{gemma-3-12b-pt} and \texttt{gemma-3-27b-pt} is not statistically significant, suggesting that most of the performance gain saturates around the 12B model.






However, when multiple models were compared, it was found that larger size did not always correspond to better syntactic accuracy. Average accuracy of \texttt{gemma-3-12b-pt}, \texttt{gemma-3-27b-pt}, \texttt{LLaMA-3-70B}, and \texttt{Alif-1.0-8B-Instruct} is 92.68,92.42, 94.73 and 93.25, respectively is comparable.
As illustrated in Figure~\ref{fig:model-comparison}, the smaller model (e.g., \texttt{gemma-3-27b-pt}) was observed to outperform larger models (\texttt{Llama-3-70B and DeepSeek-R1-Distill-Llama-70B}) on several syntactic paradigms, dative object ProperNoun, dative object pronoun, noun adjective agreement and all three paradigms of  subject-verb agreement. These results suggest that certain syntactic patterns were generalized more effectively by smaller models, potentially due to simpler architectures or reduced overfitting to training data idiosyncrasies.
 The statistical significance of this difference was confirmed through a two-tailed Wilcoxon signed-rank test.  As shown in Table~\ref{tab:pairwise_top}, no statistically significant differences ($\alpha = 0.05$) were observed among \texttt{gemma-3-12b-pt, gemma-3-27b-pt, Meta-Llama-3-70B, and Alif}, indicating that these models perform comparably in our evaluation.
\begin{table}[h]
\resizebox{\columnwidth}{!}{

\begin{tabular}{lllll}
\toprule
\rotatebox{0}{\textbf{LM1}} & \rotatebox{0}{\textbf{LM2}} & \multicolumn{1}{l}{\rotatebox{0}{\textbf{p-raw}}} & \multicolumn{1}{l}{\rotatebox{0}{\textbf{p- holm}}} & \rotatebox{0}{\textbf{Sig-holm}} \\ \hline
gemma-3-12b-pt & gemma-3-27b-pt & 0.460 & 1 & No \\
gemma-3-12b-pt & Llama-70B & 0.238 & 1 & No \\
gemma-3-12b-pt & Alif & 0.910& 1 & No \\
gemma-3-27b-pt & Llama-70B & 0.496 & 1 & No \\
gemma-3-27b-pt & Alif & 0.910 & 1 & No \\
Llama-70B & Alif & 0.195 & 1 & No \\
\hline
\end{tabular}
}
\caption{Pairwise comparisons of top-performing models using Wilcoxon signed-rank test with Holm correction. No statistically significant differences were found among the top models at $\alpha = 0.05$.}
\label{tab:pairwise_top}
\end{table}
\subsection{ Pretrained vs Instruction-Tuned Models}
A comparative evaluation was performed on the pretrained (-pt) and instruction-tuned (-it) variants of the Gemma models across four sizes: 1B, 4B, 12B, and 27B. To evaluate the statistical significance of differences between the paired model variants, a two-tailed Wilcoxon signed-rank test was applied to the accuracy scores.
We found that the pretrained variants of Gemma-3-1B, 3-4B, and 3-12B significantly outperformed their instruction-tuned counterparts. These results suggest that instruction tuning may adversely affect syntactic generalization in smaller-scale models. Fine-grained results of pretrained and instruction-tuned models against each paradigms are shown in Appendix Table  \ref{tab:pretrain-paradigm}  and \ref{tab:it-paradigm}.

\section{Conclusion}
In this study, we introduced a new benchmark 
%was introduced 
to evaluate the syntactic capabilities of multilingual language models, with a focus on under-represented and morphologically rich constructions in Urdu. The dataset was constructed following the SLING framework, using naturally formulated minimal pairs that target core syntactic phenomena. It was found that generally multilingual models performed reasonably in capturing local dependencies and agreement structures. However, reduced performance was observed on complex syntactic constructions, such as long distance agreement in subject-verb. These findings highlight the challenges that remain for language modeling in low-resource and typologically diverse languages, emphasizing the need for more linguistically informed pretraining and evaluation approaches.


\section*{Limitations}
In this study, the evaluation of linguistic phenomena in Urdu was limited primarily by the availability of annotated data. The existing Urdu Treebank contains only a small number of relevant examples for the targeted syntactic categories. To address this limitation, additional instances were extracted from unannotated Urdu corpus using carefully designed regular expressions, followed by extensive manual cleaning. Despite these efforts, the dataset size remains relatively small, especially when compared to large-scale resources like the English BLiMP or the Mandarin SLING datasets. However, we believe that our proposed dataset size is sufficient to reliably evaluate the linguistic competence of language models as shown in our results.

Furthermore, the current evaluation only covers ten linguistic categories, few more phenomenon can also be included:
\textbf{Reciprocity}: For instance,
\textit{ bace ek dosre ko ma:r rehe h~{e}}  (
    \textit{Children are hitting each other}). This construction expresses reciprocal action between subjects and acceptable in Urdu. Its variation  \textit{ bace ek dosre  ma:r rehe h~{e}}    is unacceptable.

Question constructions involving relative clauses is also not included in this dataset. For instance,  
 \textit{   jo muhabat nah~{e} rakhta vo Xuda ko nah~{e} janta:}
  (
    \textit{He who does not have love does not know God})  is acceptable but  \textit{   jo kiya: nah~{e} rakhta vo Xuda ko nah~{e} janta:} this type of sentence includes a \textit{wh}-word (\foreignlanguage{arabic}{جو}) that functions as both a relative clause marker and a question-like structure, which remains unexplored in our current dataset.

These phenomena are linguistically rich and crucial for comprehensive evaluation, and their absence points to a limitation of the current resource. With improved Urdu parsers and extended annotated corpora, such constructions could be systematically included in future datasets.


\bibliography{main}
\appendix

\section{Ethical Consideration}
We plan to open-source our dataset along with a detailed data card. The documentation will follow the templates used in the BLiMP benchmark \cite{blimp} and the HuggingFace Datasets repository \cite{huggingfaceDataset}. The dataset and accompanying code will be released via a public GitHub repository under a permissive open-source license.
\section{Computational Cost}
The computational cost of evaluating a language model (LM) on UrBLiMP varies depending on the model's architecture and size. Distributed inference libraries (e.g., accelerate, torchrun) can be used to optimize the process.

For a single NVIDIA A100 GPU with 80GB memory, the complete evaluation takes approximately 1.5 hours for decoder-only models and 10 hours for encoder-only models.

LLaMA-70B and DeepSeek-70B, due to their large size, require three A100 (80GB) GPUs for evaluation, with an estimated runtime of around 7 hours.
\section{Use of AI-Assistant}
ChatGPT was used to proofread and improve the text of this paper by correcting grammatical, spelling, and stylistic errors. 
\section{Comparison with MultiBLiMP}
\label{appendix:multiblimp}


To contextualize the difficulty of UrBLiMP and validate our findings against an existing benchmark, we compared model performance on the subject–verb agreement (SVA) phenomenon in UrBLiMP with the Urdu portion of the MultiBLiMP dataset. The Urdu portion of MultiBLiMP consists of all 550 available sentences, which exclusively test SVA, whereas UrBLiMP contains 778 SVA sentences. This enables a direct and controlled comparison between the two datasets. The results are presented in Table~\ref{tab:multiblimp_comparison}.

                                                       
\begin{table}[htbp]
\centering

\resizebox{\columnwidth}{!}{
\begin{tabular}{lrr}
\hline
\rowcolor{gray!20}
\textbf{Model} & \textbf{MultiBLiMP (SVA)} & \textbf{UrBLiMP (SVA)} \\
\hline
Gemma-3-1b-pt & 93.27 & 87.60 \\
Gemma-3-4b-pt & 94.18 & 88.91 \\
Gemma-3-12b-pt & 96.36 & 92.58 \\
Gemma-3-27b-pt & 97.09 & \textbf{94.89} \\
Gemma-3-1b-it & 76.91 & 70.44 \\
Gemma-3-4b-it & 90.18 & 82.19 \\
Gemma-3-12b-it & 95.09 & 89.99 \\
Gemma-3-27b-it & \textbf{97.27} & 92.55 \\
LLAMA-8b & 93.82 & 76.60 \\
LLAMA-70b & 96.55 & 87.13 \\
Alif-llama-8b & 94.55 & 83.13 \\
Granite-3.3-2b-base & 88.00 & 72.48 \\
Granite-3.3-8b-base & 94.36 & 64.78 \\
Granite-3.3-2b-instruct & 81.27 & 71.06 \\
Granite-3.3-8b-instruct & 88.91 & 61.07 \\
DeepSeek-R1-Distill-Llama-70B & 95.82 & 86.33 \\
bloomz & 86.00 & 84.58 \\
gpt-4o(base) &97.64 &96.70\\
gpt-4o(GP) &\textbf{98.00} &\textbf{98.06}\\
gpt-4o(GP+CoT)& 98.00&91.29\\
BERT-multilingual-base & 86.00 & 69.73 \\
T5-small & 52.36 & 50.79 \\
T5-large & 41.45 & 58.30 \\
\hline

\end{tabular}
}
\caption{Model performance comparison on Urdu Subject-Verb Agreement (SVA) phenomena. The highest score on each benchmark is highlighted in bold.}
\label{tab:multiblimp_comparison}
\end{table}

The results reveal a consistent and notable trend: for almost every model evaluated, accuracy is significantly higher on the MultiBLiMP SVA subset than on the SVA phenomena within UrBLiMP. This performance gap indicates that the subject-verb agreement constructions in UrBLiMP are more challenging, suggesting that our benchmark probes a greater range of syntactic complexity or uses more nuanced linguistic stimuli.

Furthermore, the model rankings between the two benchmarks validate the robustness of our evaluation. A key finding is that the Gamma-3-27b-it model outperformed all others, including the much larger LLAMA-70b model, on both benchmarks. This is a significant result, demonstrating that a specialized model can achieve superior grammatical understanding in Urdu compared to larger general-purpose LLMs.
\section{GPT-4 prompt templates and API metadata}
\label{appendix:prompts}

All closed-weight (API-based) evaluations were carried out with \texttt{gpt-4.0} using the OpenAI chat-completions endpoint. 
Each item consisted of a minimal pair drawn from UrBLiMP (\texttt{sentence\_good}, \texttt{sentence\_bad}).  
For every paradigm, the model was tested under three prompting regimes: \textbf{Base}, \textbf{Grammar-Prompting (GP)}, and \textbf{Grammar-Prompting + Chain-of-Thought (GP + CoT)}.  
Prompts were programmatically constructed using Python functions; random shuffling determined whether the grammatical or ungrammatical sentence appeared first (A or B).  
All calls used deterministic decoding (\texttt{temperature = 0}, \texttt{top\_p = 1}), returning one completion per prompt.

\paragraph{1. Base prompt.}
A direct two-choice formulation without any linguistic explanation:
\begin{quote}
Two sentences are given. Which one is grammatically acceptable?  

Sentence A: <sentence\_A>  

Sentence B: <sentence\_B>  

Answer exactly ``A'' or ``B'', no explanation.
\end{quote}

\paragraph{2. Grammar-Prompting (GP).}
Before presenting the pair, the model was first asked to produce a concise metalinguistic description of the relevant phenomenon, generated through:
\begin{quote}
Explain to a beginner learner the grammatical phenomenon called ``<PARADIGM>''.  
Please provide a concise metalinguistic explanation of the key rule(s) or constraint(s), but do *not* use full example sentences from the test minimal pairs.
\end{quote}

The resulting explanation text (\texttt{<EXPLANATION>}) was then prepended to the evaluation prompt:
\begin{quote}
<EXPLANATION>

Which sentence is more grammatically acceptable?  

Sentence A: <sentence\_A>  

Sentence B: <sentence\_B>  

Answer exactly ``A'' or ``B'', no explanation.
\end{quote}

\paragraph{3. Grammar-Prompting + Chain-of-Thought (GP + CoT).}
Identical to the GP prompt, but explicitly requesting step-by-step reasoning:
\begin{quote}
<EXPLANATION>

Please carefully reason step by step to answer the following question:  

Which sentence is more grammatically acceptable?  

Sentence A: <sentence\_A>  

Sentence B: <sentence\_B>  

Then give your final answer as ``A'' or ``B'' only, with no additional commentary.
\end{quote}



\onecolumn



\section{Fine-grained Results at Paradigm level}
\begin{table*}[ht!]

\resizebox{\columnwidth}{!}{
\renewcommand{\arraystretch}{1.5} % 1.5 times the default height
\begin{tabular}{llrrrrrrrrrrr}
\rowcolor{gray!20}
\multicolumn{1}{l}{\textbf{Phenomenon}} &\multicolumn{1}{l}{\textbf{Paradigm}} & \multicolumn{1}{l}{\rotatebox{80}{gemma-3-1b-pt}} & \multicolumn{1}{l}{\rotatebox{80}{gemma-3-4b-pt}} & \multicolumn{1}{l}{\rotatebox{80}{gemma-3-12b-pt}} & \multicolumn{1}{l}{\rotatebox{80}{gemma-3-27b-pt}} & \multicolumn{1}{l}{\rotatebox{80}{Llama-3-8B}} & \multicolumn{1}{l}{\rotatebox{80}{Llama-3-70B}} & \multicolumn{1}{l}{\rotatebox{80}{\shortstack{DeepSeek-R1\\Distill-Llama-70B}}} & \multicolumn{1}{l}{\rotatebox{80}{granite-3.3-2b-base}} & \multicolumn{1}{l}{\rotatebox{80}{granite-3.3-8b-base}} & \multicolumn{1}{l}{\rotatebox{80}{bloomz-7b1-p3}} & \multicolumn{1}{l}{\rotatebox{80}{\textbf{Average}}} \\ \hline\textbf{Aspect Agr} & \textbf{Aspect} & 100.00 & 100.00 & 100.00 & 100.00 & 100.00 & 100.00 & 100.00 & 99.25 & 97.51 & 99.25 & 99.60 \\ \hline
\multirow{2}{*}{\textbf{Dative-Obj}} & \textbf{Noun} & 98.75 & 98.75 & 98.75 & 100.00 & 91.25 & 94.38 & 86.25 & 87.50 & 87.50 & 99.38 & 94.25 \\ \cdashline{2-13}
 & \textbf{Pronoun} & 98.48 & 97.73 & 97.73 & 99.24 & 93.94 & 98.48 & 90.15 & 91.67 & 90.15 & 93.18 & 95.08 \\ \hline
\multirow{3}{*}{\textbf{Ergativity}} & \textbf{Verb.PEF} & 95.14 & 95.04 & 98.32 & 98.81 & 98.02 & 99.70 & 98.32 & 91.77 & 97.92 & 97.13 & 97.02 \\ \cdashline{2-13}
 & \textbf{DOM} & 82.93 & 84.92 & 88.91 & 92.02 & 97.34 & 98.45 & 97.56 & 78.94 & 85.59 & 88.91 & 89.56 \\ \cdashline{2-13}
 & \textbf{Verb-Ob Agr} & 84.00 & 90.50 & 93.00 & 94.00 & 79.00 & 91.00 & 84.50 & 75.50 & 82.50 & 86.50 & 86.05 \\ \hline
\multicolumn{1}{l}{\textbf{Exp-Subj}} & \textbf{Oblique Pronoun} & 79.01 & 84.20 & 94.07 & 84.20 & 98.27 & 95.80 & 96.30 & 77.04 & 91.11 & 67.65 & 86.77 \\ \hline
\multicolumn{1}{l}{\textbf{Honorific}} & & 73.20 & 73.86 & 84.97 & 84.31 & 93.46 & 98.04 & 99.35 & 83.66 & 90.20 & 84.31 & 86.54 \\ \hline
\multicolumn{1}{l}{\textbf{Adj-Noun Agreement}} & & 93.00 & 93.00 & 96.50 & 97.00 & 89.50 & 94.00 & 90.50 & 83.00 & 85.00 & 87.50 & 90.90 \\ 
\hline
\multirow{5}{*}{\textbf{Oblique}} & \textbf{Adjective} & 84.31 & 81.37 & 92.16 & 94.12 & 72.55 & 85.29 & 88.24 & 70.59 & 81.37 & 88.24 & 83.82 \\ \cdashline{2-13}
 & \textbf{Plural} & 88.35 & 83.50 & 89.32 & 90.29 & 85.44 & 95.15 & 88.35 & 89.32 & 92.23 & 77.67 & 87.96 \\ \cdashline{2-13}
 & \textbf{Pronoun} & 92.00 & 90.00 & 94.50 & 96.30 & 82.70 & 90.40 & 87.70 & 72.30 & 73.50 & 86.60 & 86.60 \\ \cdashline{2-13}
 & \textbf{Noun.SG.M} & 75.00 & 75.00 & 77.00 & 84.00 & 85.00 & 93.00 & 95.00 & 81.00 & 79.00 & 74.00 & 81.80 \\ \cdashline{2-13}
 & \textbf{Verb} & 97.06 & 98.04 & 98.04 & 99.02 & 100.00 & 100.00 & 99.02 & 98.04 & 98.04 & 98.04 & 98.53 \\ \cdashline{2-13}
 \hline
\textbf{Participial Relatives} &  & 79.07 & 81.06 & 85.71 & 85.38 & 90.03 & 92.69 & 88.70 & 80.07 & 82.72 & 81.06 & 84.65 \\
\hline
\multirow{3}{*}{\textbf{Subj-Verb Agre}} & \textbf{Gender} & 80.90 & 80.90 & 89.14 & 92.51 & 67.04 & 83.15 & 83.52 & 70.04 & 78.28 & 83.15 & 80.86 \\ \cdashline{2-13}
 & \textbf{Number} & 93.69 & 93.69 & 94.17 & 97.09 & 77.18 & 85.44 & 82.04 & 70.87 & 77.18 & 92.23 & 86.36 \\ \cdashline{2-13}
 & \textbf{Person} & 88.20 & 92.13 & 94.43 & 95.08 & 85.57 & 92.79 & 93.44 & 53.44 & 61.97 & 78.36 & 83.54 \\ \hline
\multicolumn{1}{l}{\textbf{Order Variation}} & & 82.18 & 86.14 & 91.09 & 91.09 & 90.10 & 94.06 & 90.10 & 85.15 & 93.07 & 85.15 & 88.81 \\
\hline
\end{tabular}
}
\caption{Fine-grained evaluation results of each syntactic paradigm using \textbf{\textit{pre-trained models}}. The final column reports the average accuracy of each paradigm across all models.}
\label{tab:pretrain-paradigm}


\end{table*}
\begin{figure*}[h!]
    \centering
    \includegraphics[width=0.99\textwidth]{images/small_vs_large.png}
    \caption{Comparison of model accuracy across linguistic phenomena, showing that the small model outperforms larger models on several linguistic phenomena.}
    \label{fig:model-comparison}
\end{figure*}

\begin{table*}[ht!]

\resizebox{\textwidth}{!}{
\renewcommand{\arraystretch}{1.5} % 1.5 times the default height
\begin{tabular}{llrrrrrrrr}
\rowcolor{gray!20}
\multicolumn{1}{l}{\textbf{Phenomenon}}  &  \multicolumn{1}{l}{\textbf{Paradigm}} & \multicolumn{1}{l}{\rotatebox{80}{gemma-3-1b-it}} & \multicolumn{1}{l}{ \rotatebox{80}{gemma-3-4b-it}} & \multicolumn{1}{l}{\rotatebox{80}{gemma-3-12b-it}} & \multicolumn{1}{l}{\rotatebox{80}{gemma-3-27b-it}} & \multicolumn{1}{l}{\rotatebox{80}{Alif}} & \multicolumn{1}{l}{\rotatebox{80}{granite-3.3-8b-it}} & \multicolumn{1}{l}{\rotatebox{80}{granite-3.3-2b-it}} & \multicolumn{1}{l}{\textbf{\rotatebox{80}{Average}}} \\ \hline
\textbf{Aspect} & \textbf{Aspect} & 99.00  & 99.75  & 100.00  & 100.00  & 100.00  & 98.75  & 99.00  & 99.50  \\ \hline
\multirow{2}{*}{\textbf{Dative-Obj}} & \textbf{Noun} & 92.50  & 94.38  & 95.62  & 95.62  & 93.12  & 87.50  & 80.00  & 91.25  \\ \cdashline{2-10}
 & \textbf{Pronoun} & 87.12  & 90.91  & 94.70  & 93.18  & 89.39  & 90.91  & 81.06  & 89.61  \\ \hline
\multicolumn{1}{l}{} & \textbf{Verb.PEF} & 76.02  & 91.77  & 95.74  & 98.71  & 98.61  & 93.95  & 87.61  & 91.77  \\ \cdashline{2-10}
\textbf{Ergativity} & \textbf{DOM} & 56.54  & 81.60  & 86.92  & 90.24  & 98.23  & 84.48  & 74.94  & 81.85  \\ \cdashline{2-10}
\multicolumn{1}{l}{} & \textbf{Verb-Obj Agr} & 77.50  & 86.50  & 92.50  & 93.00  & 84.00  & 74.00  & 78.00  & 83.64  \\ \hline
\multicolumn{1}{l}{\textbf{Experiencer-Subj}} & \textbf{Oblique Pronoun} & 62.72  & 56.54  & 76.05  & 56.05  & 98.02  & 60.74  & 92.10  & 71.75  \\ \hline
\multicolumn{1}{l}{\textbf{Honorific}} & & 55.56  & 74.51  & 79.74  & 84.97  & 98.04  & 94.77  & 88.89  & 82.35  \\ \hline
\multicolumn{1}{l}{\textbf{Adj-Noun Agreement}} & & 76.50  & 77.00  & 89.50  & 92.50  & 89.00  & 80.50  & 74.00  & 82.71  \\  \hline
\multirow{5}{*}{\textbf{Oblique}} & \textbf{Adjective} & 69.61  & 78.43  & 88.24  & 89.22  & 83.33  & 78.43  & 62.75  & 78.57  \\ \cdashline{2-10}
 & \textbf{Plural} & 66.02  & 81.55  & 78.64  & 84.47  & 90.29  & 85.44  & 81.55  & 81.14  \\ \cdashline{2-10}
 & \textbf{Pronoun} & 58.50  & 59.80  & 64.10  & 75.50  & 89.00  & 65.70  & 57.90  & 67.21  \\ \cdashline{2-10}
 & \textbf{Noun.SG.M} & 66.00  & 70.00  & 76.00  & 79.00  & 92.00  & 61.00  & 69.00  & 73.29  \\ \cdashline{2-10}
 & \textbf{Verb} & 90.20  & 94.12  & 92.16  & 97.06  & 99.02  & 99.02  & 99.02  & 95.80  \\ \hline 
\textbf{Participial Relatives} &  & 57.48  & 74.42  & 78.74  & 85.38  & 91.69  & 74.09  & 78.41  & 77.17  \\ \hline
\multirow{3}{*}{\textbf{Subj-Verb Agreement}} & \textbf{Gender} & 57.30  & 71.54  & 84.64  & 91.01  & 80.90  & 77.53  & 69.29  & 76.03  \\ \cdashline{2-10}
 & \textbf{Number} & 81.55  & 90.78  & 94.17  & 95.15  & 77.67  & 69.42  & 63.11  & 81.69  \\ \cdashline{2-10}
 & \textbf{Person} & 72.46  & 84.26  & 91.15  & 91.48  & 90.82  & 66.23  & 50.82  & 78.17  \\ \hline
\multicolumn{1}{l}{} & \textbf{Order Variation} & 63.37  & 68.32  & 78.22  & 83.17  & 97.03  & 91.09  & 76.24  & 79.63 \\ \hline  

\end{tabular}
}
\caption{Fine-grained evaluation results of each syntactic paradigm using \textbf{\textit{Instruction-tuned models}}. The final column reports the average accuracy of each paradigm across all models.}

\label{tab:it-paradigm}

\end{table*}
\clearpage
\section{Human and Model Accuracy on the Validated Subset}
\label{appendix:human_comparison}

To directly compare human and model judgments, we evaluated all models on the same subset of 1,077 sentence pairs that were manually annotated by human evaluators.  
The Table \ref{tab:human_llm_subset_appendix}  reports accuracy (\%) across linguistic phenomena for humans and representative language models.

\begin{table*}[h!]
\centering
\small


\begin{tabular}{lllllllllll}

\textbf{Model} & \rotatebox{70}{\textbf{Aspect Agr}} & \rotatebox{70}{\textbf{Dative-Obj} }& \rotatebox{70}{\textbf{Ergativity}} & \rotatebox{70}{\textbf{Exper-Sub} }&\rotatebox{70}{ \textbf{Honorific}} & \rotatebox{70}{\textbf{N-JAgr}} & \rotatebox{70}{\textbf{Oblique}} & \rotatebox{70}{\textbf{PR} }& \rotatebox{70}{\textbf{S-Verb Agr}} & \rotatebox{70}{\textbf{Order Variation} }\\ \hline
gemma-3-1b-pt & 100.00 & 98.75 & 92.50 & 82.05 & 65.00 & 95.00 & 86.50 & 80.00 & 95.15 & 77.50 \\
gemma-3-4b-pt & 100.00 & 97.50 & 94.17 & 92.31 & 72.50 & 87.50 & 85.50 & 82.50 & 93.20 & 82.50 \\
gemma-3-12b-pt & 100.00 & 97.50 & 95.00 & 94.87 & 85.00 & 92.50 & 92.50 & 87.50 & 97.09 & 87.50 \\
gemma-3-27b-pt & 100.00 & \textbf{100.00} & \textbf{98.33} & 92.31 & 77.50 & 95.00 & 94.50 & 92.50 & 97.09 & 85.00 \\
gemma-3-1b-it & 100.00 & 90.00 & 75.83 & 58.97 & 45.00 & 75.00 & 68.00 & 57.50 & 72.82 & 60.00 \\
gemma-3-4b-it & 100.00 & 95.00 & 90.83 & 58.97 & 75.00 & 72.50 & 79.00 & 67.50 & 89.32 & 65.00 \\
gemma-3-12b-it & 100.00 & 96.25 & 93.33 & 76.92 & 72.50 & 87.50 & 81.00 & 80.00 & 95.15 & 82.50 \\
gemma-3-27b-it & 100.00 & 96.25 & 95.00 & 61.54 & 87.50 & 97.50 & 85.50 & 85.00 & 95.15 & 77.50 \\
Meta-Llama-3-8B & 100.00 & 88.75 & 91.67 & \textbf{97.44} & 97.50 & 87.50 & 87.50 & 90.00 & 91.26 & 85.00 \\
Meta-Llama-3-70B & 100.00 & 96.25 & 95.83 & 94.87 & \textbf{100.00} & 92.50 & 93.00 & \textbf{100.00} & 94.17 & 92.50 \\
Alif-1.0-8B-Instruct & 100.00 & 85.00 & 94.17 & 97.44 & \textbf{100.00} & 90.00 & 90.50 & 92.50 & 94.17 & \textbf{97.50} \\
granite-3.3-2b-base & 100.00 & 88.75 & 83.33 & 71.79 & 77.50 & 77.50 & 85.00 & 75.00 & 70.87 & 82.50 \\
granite-3.3-8b-base & 97.50 & 91.25 & 89.17 & 94.87 & 90.00 & 77.50 & 83.50 & 82.50 & 78.64 & 90.00 \\
granite-3.3-2b-instruct & 100.00 & 80.00 & 79.17 & 89.74 & 82.50 & 70.00 & 73.50 & 80.00 & 66.02 & 65.00 \\
granite-3.3-8b-instruct & 97.50 & 88.75 & 83.33 & 56.41 & 92.50 & 82.50 & 77.00 & 70.00 & 75.73 & 90.00 \\
DeepSeek-R1-Distill-Llama-70B & 100.00 & 87.50 & 92.50 & 94.87 & 100.00 & 92.50 & 93.00 & 95.00 & 93.20 & 90.00 \\
gpt-4o (base) & 100.00 & 91.25 & 94.17 & 94.87 & 90.00 & \textbf{100} & 97.00 & 97.50 & \textbf{98.06} & 97.50 \\
gpt-4o(GP) & 100.00 & 92.50 & 92.50 & 97.44 & 95.00 & 100 & \textbf{98.50} & 100 & 97.09 & 97.50 \\
gpt4o(GP+CoT) & 100.00 & 95.00 & 91.67 & 97.44 & 100.00 & 97.50 & 97.00 & 97.50 & 94.17 & 100 \\
bloom-7b1 & 100.00 & 97.50 & 98.33 & 79.49 & 87.50 & 90.00 & 87.00 & 77.50 & 90.29 & 87.50 \\
bert & 75.23 & 78.75 & 78.00 & 80.33 & 75.26 & 76.04 & 80.14 & 78.79 & 76.92 & 77.35 \\
mt5-small & 53.91 & 41.25 & 55.00 & 60.25 & 56.27 & 54.03 & 57.78 & 59.26 & 59.11 & 59.02 \\
mt5-large & 59.24 & 36.25 & 50.5 & 56.06 & 57.7 & 57.93 & 61.35 & 60.93 & 62.39 & 62.80 \\
\rowcolor[HTML]{D9D9D9}
\textbf{Human Evaluation} & \textbf{100.00} & \textbf{94.27} & \textbf{97.71} & \textbf{93.89} & \textbf{94.17} & \textbf{97.29} & \textbf{95.46} & \textbf{95.63} & \textbf{95.68} & \textbf{98.75}\\
\hline
\end{tabular}
\caption{Accuracy (\%) of humans and LLMs on the human-validated subset (N ≈ 1,077).  
Grey row marks human performance.}
\label{tab:human_llm_subset_appendix}
\end{table*}




\section{UrBLiMP Examples}
\label{sec:appendix}
This appendix contains examples from each of the 19 paradigms in UrBLiMP as shown in Table \ref{tab:paradigm_examples}
\newline

\begin{table}[h!]
\resizebox{\textwidth}{!}{
\renewcommand{\arraystretch}{2.5}
\begin{tabular}{lllll}

\textbf{Phenomenon} & \textbf{Paradigm} & \textbf{N} & \textbf{Grammatical Sentence} & \textbf{Ungrammatical Sentence} \\ \hline
%%%%%%%%%%First Row %%%%%%%%%%%%%%%%%%%%%%%%%%%%%

\textbf{Aspect} &  & 400 &
 \makecell[tl]{%
      \begin{tabular}{@{}l@{}}
        \foreignlanguage{arabic}{ہاں وہ گورنر سے ملتا \underline{رہا} تھا۔} \\
        h\~{a} \quad voh \quad gavarnar \quad se \quad milta: \quad raha: \quad tha: \\
        Yes \quad he \quad governor \quad from \quad meet-\textsc{impf} \quad \textsc{prog.part} \quad was \\
        `Yes, he had been meeting the governor.' \\
     \end{tabular}
} &
\makecell[tl]{%
\begin{tabular}{@{}l@{}}
    \foreignlanguage{arabic}{ہاں وہ گورنر سے ملتا \underline{چکا} تھا۔} \\
    h\~{a} \quad voh \quad gavarnar \quad se \quad milta: \quad cuka: \quad tha: \\
    Yes \quad he \quad governor \quad from \quad meet-\textsc{impf} \quad \textsc{perf.part} \quad was \\
    `Yes, he had already met the governor.' \\
 \end{tabular}
} \\
\hline

%%%%%%%%%%%%%%%Second Row %%%%%%%%%%%%%%%
\multirow{2}{*}{\textbf{Dative Object}} 
 & Noun & 160 & 
 \makecell[tl]{%
  \begin{tabular}{@{}l@{}}
   \foreignlanguage{arabic}{پہرے‌دار نے تمام ماجرا راجا \underline{کو} سنایا۔} \\
   pahreda:r \quad ne \quad tama:m \quad ma:jra: \quad ra:ja: \quad ko \quad suna:ya: \\
   guard \quad \textsc{erg} \quad whole \quad incident \quad Raja \quad \textsc{dat} \quad narrated.\textsc{pst} \\
   `The guard narrated the whole incident to Raja.' \\
   \end{tabular}
 } 
 & 
 \makecell[tl]{%
    \begin{tabular}{@{}l@{}}
   \foreignlanguage{arabic}{پہرے‌دار نے تمام ماجرا راجا سنایا۔} \\
   pahreda:r \quad ne \quad tama:m \quad ma:jra: \quad ra:ja: \quad suna:ya: \\
   guard \quad \textsc{erg} \quad whole \quad incident \quad Raja \quad narrated.\textsc{pst} \\
   `The guard narrated the whole incident to Raja.' \\
   \end{tabular}
 } 
 \\
\cdashline{2-5}
 & Pronoun & 131 & 
 \makecell[tl]{%
  \begin{tabular}{@{}l@{}}
   \foreignlanguage{arabic}{آج مجھ \underline{کو} فرصت ہے} \\
   a:̃j \quad mujh \quad ko \quad fursat \quad hai \\
   today \quad me.\textsc{dat} \quad free.time \quad is \\
   `Today I have free time.' \\
   \end{tabular}
 }
 &
 \makecell[tl]{%
  \begin{tabular}{@{}l@{}}
   \foreignlanguage{arabic}{آج مجھ فرصت ہے} \\
   a:̃j \quad mujh \quad fursat \quad hai \\
   today \quad me \quad free.time \quad is \\
   `Today I have free time.' \\
   \end{tabular}
 } \\
\hline
\multirow{3}{*}{\textbf{Ergativity}} 
 & Perf.Verb & 1009 & 
 \makecell[tl]{%
   \begin{tabular}{@{}l@{}}
   \foreignlanguage{arabic}{پہرے‌دار نے تمام ماجرا راجا کو سنایا۔} \\
   pahreda:r \quad ne \quad tama:m \quad ma:jra: \quad ra:ja: \quad ko \quad suna:ya: \\
   guard \quad \textsc{erg} \quad whole \quad incident \quad Raja \quad \textsc{dat} \quad narrated.\textsc{pst} \\
   `The guard narrated the whole incident to Raja.' \\
   \end{tabular}
 }
 & 
 \makecell[tl]{%
   \begin{tabular}{@{}l@{}}
   \foreignlanguage{arabic}{پہرے‌دار نے تمام ماجرا راجا کو سناتا۔} \\
   pahreda:r \quad ne \quad tama:m \quad ma:jra: \quad ra:ja: \quad ko \quad suna:ta: \\
   guard \quad \textsc{erg} \quad whole \quad incident \quad Raja \quad \textsc{dat} \quad narrates.\textsc{pres} \\
   `The guard narrates the whole incident to Raja.' \\
   \end{tabular}
 } 
 \\
\cdashline{2-5}

 & DOM & 451 & 
 \makecell[tl]{%
   \begin{tabular}{@{}l@{}}
   \foreignlanguage{arabic}{آخر آپا نے منجھلی لڑکی کو آزمایا} \\
   a:khir \quad a:pa: \quad ne \quad manjhli: \quad laṛki: \quad ko \quad a:zma:ya: \\
   finally \quad sister \quad \textsc{erg} \quad middle-aged \quad girl \quad \textsc{dat} \quad tested.\textsc{pst} \\
   `Finally, sister tested the middle-aged girl.' \\
   \end{tabular}
 }
 & 
 \makecell[tl]{%
   \begin{tabular}{@{}l@{}}
   \foreignlanguage{arabic}{آخر آپا نے منجھلی لڑکی کو آزمائی} \\
   a:khir \quad a:pa: \quad ne \quad manjhli: \quad laṛki: \quad ko \quad a:zma:'i: \\
   finally \quad sister \quad \textsc{erg} \quad middle-aged \quad girl \quad \textsc{dat} \quad tests.\textsc{pres} \\
   `Finally, sister tests the middle-aged girl.' \\
   \end{tabular}
 } 
 \\
\cdashline{2-5}

 & Verb-Obj Agr & 200 & 
 \makecell[tl]{%
   \begin{tabular}{@{}l@{}}
   \foreignlanguage{arabic}{ایک کسان نے چار ایکڑ گندم بوئی} \\
   ek \quad kisa:n \quad ne \quad ca:r \quad aikR \quad gandum \quad bo'i: \\
   one \quad farmer \quad \textsc{erg} \quad four \quad acre \quad wheat \quad sowed.\textsc{pst} \\
   `One farmer sowed four acres of wheat.' \\
   \end{tabular}
 }
 & 
 \makecell[tl]{%
   \begin{tabular}{@{}l@{}}
   \foreignlanguage{arabic}{ایک کسان نے چار ایکڑ گندم بویا} \\
   ek \quad kisa:n \quad ne \quad ca:r \quad aikR \quad gandum \quad boya: \\
   one \quad farmer \quad \textsc{erg} \quad four \quad acre \quad wheat \quad sowed.\textsc{pst} \\
   `One farmer sowed four acres of wheat.' \\
   \end{tabular}
 } \\
\hline
% Sample entry (you can use similar format for others)
\textbf{Expe-Subj} & Obl.Pronoun & 405 & 
\makecell[l]{
 \begin{tabular}{@{}l@{}}
\foreignlanguage{arabic}{ مجھے یہ بلاگ پوسٹ پسند آیا۔} \\
\textit{mujhe yeh blog post pasand a:ya:} \\
\textit{to.me this blog post liking came} \\
I liked this blog post.
\end{tabular}
}
& 
\makecell[l]{\foreignlanguage{arabic}{ میں یہ بلاگ پوسٹ پسند آیا۔ } \\
\textit{maĩ yeh blog post pasand a:ya:} \\
\textit{I this blog post liking came} \\
I liked this blog post.} \\
\hline

\textbf{Honorific} & & 153 & 
\makecell[l]{\foreignlanguage{arabic}{ بہن جی مجھ سے ناراض تھیں }\\
\textit{behan ji mujh se nara:z th\~{i}} \\
\textit{sister from.me upset was (hon.)} \\
Sister was upset with me.}
& 
\makecell[l]{\foreignlanguage{arabic}{ بہن جی مجھ سے ناراض تھی} \\
\textit{behan ji mujh se nara:z thi} \\
\textit{sister from.me upset was} \\
Sister was upset with me.} \\
\hline

\textbf{N-J Agr} & & 200 & 
\makecell[l]{\foreignlanguage{arabic}{ ہر مجسمہ بیس فٹ اونچا ہے ۔} \\
\textit{har mujasma bi:s fit o:nca: hai} \\
\textit{every statue twenty feet tall is} \\
Every statue is twenty feet tall.}
& 
\makecell[l]{\foreignlanguage{arabic}{ ہر مجسمہ بیس فٹ اونچی ہے ۔} \\
\textit{har mujasma bi:s fit onci: hai} \\
\textit{every statue twenty feet tall.F is} \\
Every statue is twenty feet tall.} \\
\hline
\textbf{Participial Relatives} & & 301 &
\begin{tabular}[t]{@{}l@{}}
\foreignlanguage{arabic}{آسانی سے کھولا گیا دروازہ بند کیا جا سکتا ہے} \\
\textit{a:sani: se khola: gaya: darwaza: band kiya: ja: sakta: hai} \\
easy-\textsc{adv} from open-\textsc{ptcp} go-\textsc{ptcp.m.sg} door close do-\textsc{perf.m.sg} go can be \\
‘The door that was easily opened can be closed.’
\end{tabular}
&
\begin{tabular}[t]{@{}l@{}}
\foreignlanguage{arabic}{آسانی سے کھولتا گیا دروازہ بند کیا جا سکتا ہے} \\
\textit{a:sani: se kholta: gaya: darwaza: band kiya: ja: sakta: hai} \\
easy-\textsc{adv} from open-\textsc{ipfv.m.sg} go-\textsc{ptcp.m.sg} door close do-\textsc{perf.m.sg} go can be \\
‘The door that kept opening easily can be closed.’
\end{tabular} \\
\hline

\multirow{3}{*}{\textbf{Subj-Verb Agr}} 
& Gender & 267 &
\begin{tabular}[t]{@{}l@{}}
\foreignlanguage{arabic}{اور اس کا باپ تو شاید پاگل ہو جاتا ۔} \\
\textit{aur us ka ba:p to shayad pagal ho jata:} \\
and 3\textsc{sg.m} \textsc{gen.m} father \textsc{focus} perhaps mad become \textsc{cond.m.sg} \\
‘And his father might have gone mad.’
\end{tabular}
&
\begin{tabular}[t]{@{}l@{}}
\foreignlanguage{arabic}{اور اس کا باپ تو شاید پاگل ہو جاتی ۔} \\
\textit{aur us ka ba:p to shayad pagal ho jati:} \\
and 3\textsc{sg.m} \textsc{gen.m} father \textsc{focus} perhaps mad become \textsc{cond.f.sg} \\
‘And his father (incorrectly) might have gone mad (feminine verb).’
\end{tabular}
\\ \cdashline{2-5}

& Number & 206 &
\begin{tabular}[t]{@{}l@{}}
\foreignlanguage{arabic}{مٹکا تلاش کیا تو وہ غائب تھا} \\
\textit{matka: tala:sh kiya: to vo ghaib tha:} \\
pot search did then he absent was.\textsc{sg} \\
‘When I searched for the pot, it was gone.’
\end{tabular}
&
\begin{tabular}[t]{@{}l@{}}
\foreignlanguage{arabic}{مٹکا تلاش کیا تو وہ غائب تھے} \\
\textit{matka: tala:sh kiya: to vo ghaib the} \\
pot search did then he absent was.\textsc{pl} \\
‘When I searched for the pot, they were gone.’ (number disagreement)
\end{tabular}
\\ \cdashline{2-5}

& Person & 305 &
\begin{tabular}[t]{@{}l@{}}
\foreignlanguage{arabic}{میں اسے شوق سے پڑھتا ہوں ۔} \\
\textit{mEn ise shauq se parhta: ho:n} \\
I him passion with read.\textsc{1sg.m} be.\textsc{1sg} \\
‘I read it with passion.’
\end{tabular}
&
\begin{tabular}[t]{@{}l@{}}
\foreignlanguage{arabic}{میں اسے شوق سے پڑھتا ہو ۔} \\
\textit{mEn ise shauq se parhta: ho} \\
I him passion with read.\textsc{1sg.m} be.\textsc{3sg} \\
‘I read it with passion.’ (person mismatch)
\end{tabular}
\\ \hline

\textbf{Order Variation} &  & 101 &
\begin{tabular}[t]{@{}l@{}}
\foreignlanguage{arabic}{جو بات ہے وہ بولو۔} \\
\textit{jo ba:t hai vo bolo} \\
what matter is that say.\textsc{imp.pl} \\
‘Say what the matter is.’
\end{tabular}
&
\begin{tabular}[t]{@{}l@{}}
\foreignlanguage{arabic}{جو ہے وہ بولو بات ۔} \\
\textit{jo hai vo bolo ba:t} \\
what is that say matter \\
‘Say what is that matter.’ (incorrect order)
\end{tabular}
\\ \hline


\end{tabular}

}
\caption{One representative example of grammatical and ungrammatical sentence pairs is shown for each syntactic paradigm. The minimal difference between each pair is underlined. N denotes the total number of sentence pairs included in each paradigm.}
\label{tab:paradigm_examples}
\end{table}



\end{document}
